\section{Study plan, checklist, and reading list}
\label{ch:plan}

\subsection{The four-hour version}

Three hours forty-five minutes for the day before the Koehn meeting, assuming
linear algebra from coursework and no transformers or analysis of variance.
Chapter~\ref{ch:lfs} precedes the statistics on purpose.

\begin{description}[leftmargin=3.2cm, style=nextline, itemsep=3pt]
\item[0:00 to 0:15, vectors] Chapter~\ref{ch:ml}, why a hidden state is a vector. Exercise: cosine similarity and a first look at anisotropy.
\item[0:15 to 0:50, the model] Chapter~\ref{ch:llm}: tokenization and fertility, embeddings and the residual stream, attention read once. Figure~\ref{fig:block}. Three facts: layer $0$ is the embedding output, every layer shares one $D$-dimensional space, pooling is part of the measurement.
\item[0:50 to 1:05, translation data] Chapter~\ref{ch:mt}: parallel data, translationese, the MEXA paragraph. Figure~\ref{fig:g-r6}, rows read as resource tiers.
\item[1:05 to 1:55, LFS] Chapter~\ref{ch:lfs} in full. Rederive the whiteboard table until $1.5/17.5 = 0.086$ comes unaided; type the exercise ``LFS from scratch'' and check its three printed lines. Figure~\ref{fig:grid}.
\item[1:55 to 2:30, statistics] Chapter~\ref{ch:stats}: the decomposition and bootstrap subsections. Exercise: verify the decomposition and the null by simulation. Then Figure~\ref{fig:bootstrap}: say why models, not rows, are the unit.
\item[2:30 to 3:00, results] Chapter~\ref{ch:figures}, the two paper figures. Figure~\ref{fig:g-p1}: axes, one line, the null line at $0.298$. Then Figure~\ref{fig:g-p2} and the ten numbers in \texttt{docs/KOEHN\_MEETING\_BRIEF\_2026-09-05.md}.
\item[3:00 to 3:20, limits] Chapter~\ref{ch:lfs}: the boundary subsection and what LFS relates to. Figure~\ref{fig:collapse_trained}. Then the MEXA, pooling, and ANOVA answers in Chapter~\ref{ch:defense}.
\item[3:20 to 3:30, aloud] The ten-second, thirty-second, and two-minute versions, then the whiteboard example cold, standing.
\item[3:30 to 3:45, self-test] The ten questions below, answers written first.
\end{description}

\textbf{Skipped relative to the eight-hour plan, and why that is safe.}
\begin{itemize}[itemsep=1pt]
\item Chapter~\ref{ch:ml} regression, gradient descent, softmax, XOR, PCA: LFS trains nothing; in the meeting least squares is the one phrase ``after controlling for exposure, family, script, and fertility''.
\item Chapter~\ref{ch:llm} attention exercise, MLP, normalization, training loop, logit lens: LFS reads the residual stream and never generates; Koehn's recorded questions concern pooling and the shape.
\item Chapter~\ref{ch:mt} alignment history, the ``think in English'' list, the Aim~2 options: the brief scripts the one answer that needs them.
\item Chapter~\ref{ch:stats} residualization and cluster-bootstrap exercises, REML, multiple comparisons: read the figures; the intervals are already computed.
\item The eighteen report figures and the misalignment decomposition: the meeting opens on paper Figure~1 and stays there.
\item The full question bank: the self-test replaces it; return to the bank if time remains.
\end{itemize}

\textbf{Self-test, fifteen minutes.}
\begin{enumerate}[itemsep=1pt]
\item LFS and its parts? $\SSl/(\SSl+\SSc)$: weighted spread of language means over that plus the sentence-mean spread; residual reported, not included.
\item The whiteboard value? $1.5/17.5 = 0.086$; languages differ by a constant offset, sentences by far more.
\item Why standardize jointly? Per-language standardizing subtracts the language means being measured.
\item Pure-noise value here? $(L-1)/(L+N-2) = 0.298$ at 128 by 300; compare only at a fixed grid.
\item Layer $0$, and why layers compare? Embedding output; the residual stream is a running sum, so all layers share one space.
\item Dip depth range, tracking what? $0.039$ (BLOOM-1.7B) to $0.336$ (Qwen3-8B); recipe more than size.
\item The behavioral result, with unit? Spearman $0.508$, model bootstrap $[0.19, 0.71]$, 19 models, 13 families, after controls.
\item What LFS cannot see? Uniform shrinkage: norm 451 to 42, concept variance 13{,}791 to 699, LFS $0.445$ to $0.493$; hence raw components beside the ratio.
\item MEXA in one line? Is the English translation a sentence's nearest neighbour among $N$ candidates; a published reference measure for a different question.
\item Why resample models, not rows? Rows within a model are correlated and the claim is about models.
\end{enumerate}

\subsection{Three weeks to the paper, then the long arc}

\begin{description}[leftmargin=2.9cm, style=nextline, itemsep=3pt]
\item[Days 1 to 3: own the instrument] Chapter~\ref{ch:lfs}. Rederive the worked example and the null by hand; type the \texttt{lfs} exercise; open \texttt{results/grid/Qwen3-8B-Base/metrics.json} and explain the profile layer by layer. Read \texttt{reports/lfs\_math.pdf} and the paper draft Sections 3 and 5.
\item[Days 3 to 6: the transformer] Chapter~\ref{ch:llm} and Chapter~\ref{ch:ml} Sections 1.4 to 1.7. Do the attention and XOR exercises. Read Jurafsky and Martin's transformer and LLM chapters; skim Vaswani et al.\ (2017); read the residual-stream section of Elhage et al.\ (2021); read nostalgebraist's logit-lens post and Timkey and van Schijndel (2021).
\item[Days 6 to 10: statistics] Chapter~\ref{ch:stats}. Do the decomposition, null, residualization, and cluster-bootstrap exercises. Read one ANOVA chapter (any introductory text), Efron and Tibshirani chapters 6 and 7, and Gelman and Hill on fixed effects and clustered data.
\item[Days 8 to 14: Koehn's literature] Chapter~\ref{ch:mt}. Read Koehn (2020) on the transformer, multilingual models, and evaluation; Federmann et al.\ (2022) NTREX; Conneau et al.\ (2018); Marchisio et al.\ (2020 to 2022); Kargaran et al.\ (2025) MEXA in full; Bandarkar et al.\ (2024) Belebele; Ahia et al.\ (2023); Briakou et al.\ (2023).
\item[Days 10 to 16: the debate] Wendler et al.\ (2024); Shani and Basirat (2025); Tezuka and Inoue (2025); Zhao et al.\ (2024); Bafna et al.\ (2025); Dumas et al.\ (2025); Kornblith et al.\ (2019); Chang, Tu, and Bergen (2022). For each, one sentence on its claim and one on what LFS adds or contradicts.
\item[Ongoing: foundations] Chapter~\ref{ch:ml} in full; Goodfellow, Bengio, and Courville chapters 2, 5, 6, 8, 10; Murphy's \emph{Probabilistic Machine Learning: An Introduction} on linear models and inference; a linear algebra refresher (Strang, or the \emph{Deep Learning} book chapter 2).
\end{description}

\subsection{Practice}

Say the ten-second, thirty-second, and two-minute versions aloud once a day.
Once a week, have someone read Chapter~\ref{ch:defense} questions in random
order. Before the Koehn meeting, redo the whiteboard example cold. Every
number you say should be one you can point to in the assessment document or
the claims ledger.

\subsection{The concept checklist}

Tick each when you can explain it to a first-year student in two minutes
without notes. Items marked $\star$ are the must-know set for defending
LFS; the rest are what makes you credible and lets you answer the follow-up.

\paragraph{Why learn at all.}
$\star$ rules versus examples; $\star$ parameters, loss, optimization;
$\star$ supervised, unsupervised, self-supervised, reinforcement;
$\star$ where pretraining, post-training, and measurement sit.

\paragraph{Machine learning basics.}
model, parameters, loss; $\star$ vectors, dot product, norm, cosine; matrix
as many dot products; $\star$ linear regression, least squares, residuals,
$R^2$; $\star$ ridge; gradient, gradient descent, learning rate, batches,
Adam; sigmoid, softmax, logits, cross-entropy, negative log-likelihood,
perplexity; $\star$ train/validation/test, overfitting; regularization,
early stopping, cross-validation; MLP, nonlinearity, backpropagation,
autograd; embeddings; PCA, SVD, eigenvalues, effective rank, participation
ratio; epochs, seeds, checkpoints, pretraining, fine-tuning, inference.

\paragraph{Language models.}
$\star$ next-token prediction; $\star$ tokenization, BPE, vocabulary,
fertility; embedding matrix; positional information (rotary); $\star$
residual stream; $\star$ hidden state per layer, layer 0; attention (query,
key, value, scores, causal mask, heads); MLP block; LayerNorm/RMSNorm,
pre-norm; $\star$ massive activations, anisotropy; unembedding, logits, tied
embeddings; teacher forcing; perplexity; training loop, Adam, schedules;
$\star$ base versus instruct, SFT, RLHF/DPO; greedy versus sampling,
temperature; $\star$ fp16/bf16/fp32; $\star$ mean versus last-token pooling;
$\star$ logit lens and its validity limit; linear probes; activation
patching; CKA; scaling versus recipe.

\paragraph{Machine translation and multilingual NLP.}
SMT versus NMT versus LLM; $\star$ parallel and n-way parallel corpora;
sentence alignment; $\star$ NTREX, FLORES, WMT; documents and purged splits;
$\star$ translationese and direction; BLEU, chrF, COMET; contamination;
chance floors; $\star$ resource tiers, families, scripts; curse of
multilinguality; cross-lingual transfer; incidental bilingualism; tokenizer
fairness; cross-lingual word-embedding alignment, Procrustes, MUSE, hubness,
CSLS; sentence embeddings, LASER/LaBSE; $\star$ MEXA, AaR, Information
Parity; $\star$ the ``think in English'' debate and its papers; Aim 2
methods (parallel data, code-switching, alignment losses, adversarial,
preference); Belebele, INCLUDE, Global-MMLU, MGSM.

\paragraph{Statistics of measurement.}
$\star$ variance and sums of squares; $\star$ one-way and two-way ANOVA;
$\star$ balance and orthogonality; $\star$ degrees of freedom; mean squares;
expected mean squares; variance components; REML; $\star$ standardization
and invariances; $\star$ Pearson versus Spearman; $\star$ confounding;
$\star$ residualization; $\star$ fixed effects; logit transform above a
floor; cluster-robust errors; $\star$ bootstrap; $\star$ cluster and crossed
bootstrap; $\star$ unit of generalization; effective sample size;
leave-one-out; Benjamini--Hochberg; $\star$ preregistration; minimum
detectable effect; construct and predictive validity; reliability,
ceilings, Spearman--Brown; $\star$ Goodhart's law.

\paragraph{Linear algebra and interpretability tools.}
projection; orthogonal matrices and rotations; Procrustes; PCA/SVD;
whitening; effective rank; Gini; ridge regression; linear probes; activation
patching; CKA; logit lens.

\paragraph{LFS itself.}
balanced grid; fp32 mean pooling; joint z-scoring; grand, language, sentence
means; $\SSl$, $\SSc$, $\SSr$ and the multipliers; the ratio; the residual's
role; the decomposition proof; boundedness; scale invariance; rotation
invariance; the null and its three consequences; near size-independence; dip
depth, interior minimum, endpoint recovery; values above 0.5; additive
boundary and the misalignment decomposition; variance share versus separability;
collapse blind spot and the profile fix; content transfer versus benchmarks;
observation versus intervention; the family and the no-averaging rule; MEXA
as a different question.

\subsection{Reading list with what to take from each}

\begin{description}[leftmargin=0pt, style=nextline, itemsep=2pt]
\item[Jurafsky and Martin, \emph{Speech and Language Processing}, 3rd ed.\ draft (free online)] transformers, LLMs, MT chapters: the vocabulary and the standard pictures.
\item[Koehn, \emph{Neural Machine Translation} (2020)] the transformer from an MT perspective, multilingual models, evaluation. Read for framing as much as content.
\item[Koehn, \emph{Statistical Machine Translation} (2010)] chapters on parallel data and evaluation; skim for history.
\item[Goodfellow, Bengio, Courville, \emph{Deep Learning} (free online)] ch.\ 2 linear algebra, 5 ML basics, 6 MLPs, 8 optimization, 10 sequence models.
\item[Murphy, \emph{Probabilistic Machine Learning: An Introduction} (free online)] linear models, logistic regression, inference basics.
\item[Efron and Tibshirani, \emph{An Introduction to the Bootstrap}] chapters 6 and 7; then any short note on cluster bootstrap.
\item[Gelman and Hill, \emph{Data Analysis Using Regression and Multilevel/Hierarchical Models}] fixed effects, clustered data, why the unit matters.
\item[Searle, Casella, McCulloch, \emph{Variance Components}] the balanced two-way chapter, for REML.
\item[Vaswani et al.\ 2017] attention is all you need. Read Sections 3 and 4.
\item[Elhage et al.\ 2021, \emph{A Mathematical Framework for Transformer Circuits}] the residual-stream section only.
\item[nostalgebraist 2020] the logit lens post; then Wendler et al.\ 2024.
\item[Timkey and van Schijndel 2021; Sun et al.\ 2024] rogue dimensions and massive activations.
\item[Kargaran et al.\ 2025 (MEXA)] read fully; know the exact scoring rule and the reported correlations.
\item[Conneau et al.\ 2018; Artetxe et al.\ 2018; Marchisio et al.\ 2020--2022] cross-lingual embedding alignment and Koehn's group's contributions.
\item[Shani and Basirat 2025; Tezuka and Inoue 2025; Zhao et al.\ 2024; Bafna et al.\ 2025; Dumas et al.\ 2025] the layerwise-sharing debate.
\item[Kornblith et al.\ 2019 (CKA); Chang, Tu, Bergen 2022] representational similarity and multilingual geometry.
\item[Federmann et al.\ 2022 (NTREX); Bandarkar et al.\ 2024 (Belebele); Romanou et al.\ 2025 (INCLUDE)] the datasets.
\item[Ahia et al.\ 2023; Petrov et al.\ 2023] tokenizer fertility and fairness.
\item[Briakou, Cherry, Foster 2023] incidental bilingualism drives translation ability.
\item[Nosek et al.\ 2018; Strathern 1997] preregistration; Goodhart's law.
\end{description}

\subsection{Project documents to know by location}

\begin{description}[leftmargin=0pt, style=nextline, itemsep=1pt]
\item[\texttt{docs/AIM1\_ASSESSMENT\_FOR\_ICLR\_2026-09-04.md}] sub-aim verdicts, every number with its artifact.
\item[\texttt{paper/iclr2027/main.tex}, \texttt{claims\_ledger.csv}] the paper and the map from numbers to artifacts and jobs.
\item[\texttt{docs/AIM1\_NEXT\_EXPERIMENTS\_2026-09-05.md}] the candidate experiments, costs, and frozen predictions.
\item[\texttt{code/pilot\_metrics.py}] the canonical estimator.
\item[\texttt{docs/DISCREPANCY\_LEDGER.md}] every correction, with mechanism. Reviewers reward it when you volunteer it.
\item[\texttt{proposal/NSF\_Proposal\_on\_Multilingual\_LLMs.pdf}] Section 3.1, pages 6 to 8: the questions this work answers.
\end{description}
